% !TeX program = pdflatex
% !TeX encoding = UTF-8

\documentclass[12pt,a4paper]{report}

% ===== Packages =====
\usepackage[utf8]{inputenc}
\usepackage[T1]{fontenc}
\usepackage[english]{babel}
\usepackage{geometry}
\geometry{margin=2.5cm}
\usepackage{graphicx}
\usepackage{hyperref}
\usepackage{amsmath}
\usepackage{amssymb}
\usepackage{booktabs}
\usepackage{listings}
\usepackage{xcolor}
\usepackage{natbib}

% ===== Code listings =====
\definecolor{codegreen}{rgb}{0,0.6,0}
\definecolor{codegray}{rgb}{0.5,0.5,0.5}
\definecolor{backcolour}{rgb}{0.95,0.95,0.92}

\lstdefinestyle{mystyle}{
    backgroundcolor=\color{backcolour},
    commentstyle=\color{codegreen},
    keywordstyle=\color{magenta},
    basicstyle=\ttfamily\footnotesize,
    breaklines=true,
    captionpos=b,
    numbers=left,
    numbersep=5pt,
    tabsize=2
}
\lstset{style=mystyle}

% ===== Document =====
\title{\bfseries Species Distribution Modelling with Containerized HPC}
\author{Author Name}
\date{\today}

\begin{document}

% ===== Title page =====
\begin{titlepage}
    \centering
    \vspace*{3cm}
    
    {\Huge\bfseries Species Distribution Modelling\\
    with Containerized High-Performance Computing}
    
    \vspace{2cm}
    {\Large Master's Thesis}
    
    \vspace{3cm}
    {\large Author Name}
    
    \vspace{2cm}
    {\large Department of [Department]\\
    [University]\\
    [Year]}
    
    \vfill
    {\small \today}
\end{titlepage}

% ===== Abstract =====
\chapter*{Abstract}
This thesis develops a containerized workflow for species distribution modelling (SDM) using the biomod2 R package. We implement ensemble modelling with multiple algorithms (GLM, GBM, ANN, FDA, MAXNET) to predict Alpine species distributions under current and future climate scenarios. The entire workflow is containerized using Apptainer and deployed on the Cineca Leonardo HPC cluster using SLURM scheduling. Models are trained on 62,668 European species occurrence records and evaluated with TSS, AUCroc, and KAPPA metrics. Results show ensemble methods improve predictive accuracy, and containerization ensures reproducibility across computing environments.

\vspace{0.5cm}
\noindent\textbf{Keywords:} Species Distribution Modelling, biomod2, Ensemble Methods, Climate Change, HPC, Apptainer, Cineca Leonardo

\tableofcontents

% ===== Chapter 1: Introduction =====
\chapter{Introduction}

\section{Context}

Species distribution modelling (SDM) predicts where species occur based on environmental variables. As climate changes, understanding how species distributions shift is essential for conservation planning.

This project applies SDM to Alpine species using the biomod2 R package. We combine multiple modelling algorithms to improve predictions, then containerize the workflow for reproducible execution on HPC systems.

\section{Objectives}

\begin{enumerate}
    \item Implement ensemble SDM using biomod2 with five algorithms (GLM, GBM, ANN, FDA, MAXNET)
    \item Develop a containerized workflow ensuring reproducibility
    \item Deploy models on Cineca Leonardo HPC cluster
    \item Project Alpine species distributions under current and future climate scenarios
    \item Evaluate model performance using multiple metrics
\end{enumerate}

\section{Thesis Structure}

Chapter 2 reviews SDM concepts and ensemble methods. Chapter 3 describes the data and environmental variables. Chapter 4 outlines the methodology. Chapter 5 details containerization and HPC deployment. Chapter 6 presents results. Chapter 7 concludes.

% ===== Chapter 2: Background =====
\chapter{Species Distribution Modelling}

\section{Overview}

Species distribution models relate species occurrence records to environmental predictors (climate, topography, soil) to predict geographic distributions. SDM is based on the ecological niche concept: each species occupies a region of environmental space defined by climatic and biotic factors.

\section{Algorithms}

We implement five modelling algorithms:

\begin{itemize}
    \item \textbf{GLM (Generalized Linear Model)}: Logistic regression for presence-absence data
    \item \textbf{GBM (Generalized Boosting Model)}: Ensemble of decision trees using gradient boosting
    \item \textbf{ANN (Artificial Neural Network)}: Non-linear relationships via interconnected layers
    \item \textbf{FDA (Flexible Discriminant Analysis)}: Non-linear discrimination using splines
    \item \textbf{MAXNET (Maximum Entropy)}: Probability distribution matching presence patterns
\end{itemize}

\section{Ensemble Modelling}

No single algorithm performs best across all species and regions. Ensemble modelling combines predictions from multiple algorithms to improve accuracy and reduce algorithm-specific biases. We use unweighted averaging (EMmean) to combine models exceeding AUCroc > 0.6.

\section{Evaluation Metrics}

Model performance is assessed using:
\begin{itemize}
    \item \textbf{TSS (True Skill Statistic)}: Sensitivity + Specificity - 1. Range: -1 to +1, where +1 is perfect.
    \item \textbf{AUCroc}: Area under ROC curve. Range: 0 to 1, where 1 is perfect discrimination.
    \item \textbf{KAPPA}: Agreement between predicted and observed, corrected for chance.
\end{itemize}

% ===== Chapter 3: Data and Methods =====
\chapter{Data and Methodology}

\section{Study Data}

Species occurrence data: 62,668 presence records across multiple species from European locations.

Environmental predictors (1 km resolution):
\begin{itemize}
    \item PC1clim, PC2clim: Principal components from bioclimatic variables (PCA-reduced)
    \item tri: Topographic roughness index from digital elevation models
    \item PC1soil, PC2soil: Principal components from soil properties
\end{itemize}

Models are calibrated on European data and projected to the Alpine region.

\section{Modelling Workflow}

The workflow consists of six steps:

\begin{enumerate}
    \item \textbf{Data formatting}: Combine species occurrences and environmental data; generate 50 pseudo-absences per species
    \item \textbf{Model calibration}: Fit GLM, GBM, ANN, FDA, MAXNET with 70\% training / 30\% testing split
    \item \textbf{Ensemble creation}: Combine models with AUCroc > 0.6 using unweighted averaging
    \item \textbf{Current projection}: Project models onto current Alpine environment
    \item \textbf{Future projections}: Project onto two climate scenarios (SSP3-7.0 moderate, SSP5-8.5 high emissions)
    \item \textbf{Evaluation}: Calculate TSS, AUCroc, KAPPA for all models
\end{enumerate}

\section{Implementation}

The analysis uses the \texttt{biomod2} R package with supporting libraries: terra (raster manipulation), gbm, mda, doParallel (parallel computing).

Key code structure:

\begin{lstlisting}[language=R]
# Data formatting
myBiomodData <- BIOMOD_FormatingData(
  resp.var = presence_vector,
  expl.var = environmental_rasters,
  PA.nb.absences = 50,
  PA.strategy = "random"
)

# Model calibration
myBiomodModelOut <- BIOMOD_Modeling(
  myBiomodData,
  models = c("GLM", "GBM", "ANN", "FDA", "MAXNET"),
  CV.strategy = "random",
  CV.nb.rep = 1,
  CV.perc = 0.7
)

# Ensemble
myBiomodEM <- BIOMOD_EnsembleModeling(
  bm.mod = myBiomodModelOut,
  em.algo = "EMmean",
  metric.select.thresh = 0.6
)

# Projections
myBiomodProj <- BIOMOD_Projection(
  bm.mod = myBiomodModelOut,
  new.env = current_env
)
\end{lstlisting}

% ===== Chapter 4: Containerization =====
\chapter{Containerization and HPC Deployment}

\section{Container Design}

The entire workflow is containerized using Apptainer (formerly Singularity), ensuring consistent execution across different computing environments. The container definition specifies:

\begin{itemize}
    \item \textbf{Base image}: rocker/tidyverse (R with tidyverse packages)
    \item \textbf{System libraries}: GDAL, PROJ, GEOS, udunits2 (for geospatial processing)
    \item \textbf{R packages}: biomod2, terra, raster, gbm, mda, doParallel
\end{itemize}

\begin{lstlisting}[language=bash]
Bootstrap: docker
From: rocker/tidyverse:latest

%post
    apt-get update && apt-get install -y \
        libgdal-dev libproj-dev libgeos-dev \
        libudunits2-dev libsodium-dev
    
    install2.r biomod2 terra raster gbm mda \
        Hmisc plyr doParallel

%runscript
    exec Rscript /workspace/simulation.r
\end{lstlisting}

\section{HPC Deployment}

The container is deployed on Cineca Leonardo using SLURM job scheduling:

\begin{lstlisting}[language=bash]
#!/bin/bash
#SBATCH --job-name=biomod_sdm
#SBATCH --time=24:00:00
#SBATCH --nodes=1
#SBATCH --ntasks-per-node=1
#SBATCH --cpus-per-task=10
#SBATCH --partition=boost

module load apptainer
apptainer run \
    --bind $WORK:/work,$CINECA_SCRATCH:/scratch \
    container.sif
\end{lstlisting}

The job allocates 10 CPU cores for parallel model calibration, with 24-hour walltime on the boost partition.

\section{Advantages}

\begin{itemize}
    \item \textbf{Reproducibility}: Identical results across local, HPC, and cloud environments
    \item \textbf{Portability}: Single .sif file contains all dependencies
    \item \textbf{Scalability}: SLURM integration enables dynamic resource allocation
    \item \textbf{Version control}: Container definition tracked in Git
\end{itemize}

% ===== Chapter 5: Results =====
\chapter{Results}

\section{Model Performance}

[Insert evaluation results table comparing individual and ensemble model metrics]

Ensemble models show improved TSS and AUCroc compared to individual algorithms, consistent with literature on ensemble methods. Only models with AUCroc > 0.6 were retained in the ensemble.

\section{Current Distributions}

[Insert map of current ensemble-projected species suitability]

Models identify suitable Alpine habitat constrained by climatic gradients and topography.

\section{Future Climate Scenarios}

\subsection{SSP3-7.0 (Moderate emissions)}
[Insert projection map showing moderate climate change impacts]

\subsection{SSP5-8.5 (High emissions)}
[Insert projection map showing high climate change impacts]

Comparing current and future projections reveals upward elevational shifts and potential habitat loss under both scenarios, with more severe changes under SSP5-8.5.

\section{Computational Performance}

Deployment on Cineca Leonardo reduced execution time significantly compared to local runs, enabling full processing of all 62,668 occurrence records across multiple species and future scenarios.

% ===== Chapter 6: Discussion and Conclusion =====
\chapter{Discussion and Conclusion}

\section{Key Findings}

This thesis successfully developed a containerized ensemble SDM workflow. Main contributions:

\begin{enumerate}
    \item Implemented five-algorithm ensemble using biomod2, improving predictive accuracy
    \item Containerized entire workflow with Apptainer for reproducible execution
    \item Deployed on HPC with SLURM integration for scalable processing
    \item Projected Alpine species under current and future climate scenarios
\end{enumerate}

\section{Limitations}

\begin{itemize}
    \item Models treat species independently, ignoring biotic interactions
    \item Pseudo-absence generation may introduce sampling bias
    \item Projections assume unlimited dispersal capacity
    \item Clamping masks identify extrapolation risk in novel climates
\end{itemize}

\section{Future Work}

\begin{itemize}
    \item GPU acceleration for faster model training
    \item Dynamic ensemble weighting schemes
    \item Integration of citizen science data
    \item Real-time biodiversity monitoring systems
\end{itemize}

\section{Conclusion}

This thesis demonstrates that ensemble SDM combined with containerization and HPC provides a robust, reproducible framework for biodiversity assessment under climate change. The modular workflow is extensible to other species, regions, and scenarios. As climate impacts on biodiversity intensify, such scalable, open approaches are essential for evidence-based conservation planning.

% ===== Bibliography =====
\begin{thebibliography}{99}

\bibitem{allouche2006}
Allouche, O., Tsoar, A., and Kadmon, R. (2006).
Assessing the accuracy of species distribution models: prevalence, kappa and the true skill statistic (TSS).
\emph{Journal of Applied Ecology}, 43(6):1223--1232.

\bibitem{araújo2005}
Araújo, M.B., Whittaker, R.J., Ladle, R.J., and Erhard, M. (2005).
Reducing uncertainty in projections of extinction risk from climate change.
\emph{Global Ecology and Biogeography}, 14(6):529--538.

\bibitem{elith2006}
Elith, J., Graham, C.H., Anderson, R.P., et al. (2006).
Novel methods improve prediction of species' distributions from occurrence data.
\emph{Ecography}, 29(2):129--151.

\bibitem{friedman2001}
Friedman, J.H. (2001).
Greedy function approximation: A gradient boosting machine.
\emph{Annals of Statistics}, 29(5):1189--1232.

\bibitem{guisan2005}
Guisan, A., Thuiller, W., and Zimmermann, N.E. (2005).
Predictive habitat suitability models in ecology.
\emph{Ecology Letters}, 8:993--1009.

\bibitem{hijmans2023}
Hijmans, R.J. (2023).
terra: Spatial Data Analysis.
R package version 1.7-46.

\bibitem{kurtzer2017}
Kurtzer, G.M., Sochat, V., and Bauer, M.W. (2017).
Singularity: Scientific containers for mobility of compute.
\emph{PLOS ONE}, 12(5):e0177459.

\bibitem{marmion2009}
Marmion, M., Parviainen, M., Luoto, M., Thuiller, W., and Heikkinen, R.K. (2009).
Evaluation of consensus methods in predictive species distribution modelling.
\emph{Diversity and Distributions}, 15(1):59--69.

\bibitem{phillips2006}
Phillips, S.J., Anderson, R.P., and Schapire, R.E. (2006).
Maximum entropy modeling of species geographic distributions.
\emph{Ecological Modelling}, 190(3-4):231--259.

\bibitem{thuiller2009}
Thuiller, W., Lafourcade, B., Engler, R., et al. (2009).
BIOMOD -- a platform for ensemble forecasting of species distributions.
\emph{Ecography}, 32(2):123--132.

\end{thebibliography}

\end{document}
